% !PRINTED 746
\weberpage{746}
%<<CONT
\gesperrt{one-to-one relation with the natural sequence of numbers or with a part of it}\footnote{\gesperrt{Cantor}, l.\,c. The concept of enumerable sets coincides with the concept of \gesperrt{simply infinite systems}, which \gesperrt{Dedekind} defined without the \gesperrt{assumption of the system of numbers}. (\gesperrt{Dedekind}, Was sind und was sollen die Zahlen? \S.\,6. Braunschweig 1887. Zweite unveränderte Auflage 1893.)}.

Accordingly every finite set is enumerable, and in the enumeration the natural sequence of numbers is used only up to a certain highest number. In what follows we shall consider chiefly infinite sets.

We may also give the definition of an infinite enumerable set the following expression: it is a set in which to every element a definite number of the natural sequence of numbers can be assigned as a name, in such a way that every number of the sequence of numbers is thereby used as well; that is, a set in which there is a first, a second, a third \ldots a hundredth \ldots element.

Finally we may also say that an infinite enumerable set is one that can be ordered in a sequence in such a way that a \gesperrt{first} element is present, and that upon every element there follows a \gesperrt{definite other element of the set}, and that every element, with the exception of the first, is preceded by a \gesperrt{definite other element}. Such a sequence we may call a countable arrangement.

It is clear that an enumerable set is enumerable not merely in one way, but in infinitely many different ways.

Besides the natural sequence of numbers itself, which is certainly enumerable, we may adduce as a second example the system of rational positive proper fractions, which, among other ways, can be enumerated in the following manner:
\[
\frac{1}{2}, \frac{1}{3}, \frac{2}{3}, \frac{1}{4}, \frac{3}{4}, \frac{1}{5}, \frac{2}{5}, \frac{3}{5}, \frac{4}{5}, \frac{1}{6}, \frac{5}{6}, \ldots
\]
i.\,e. in such a way that every larger denominator follows the smaller denominator, and that the fractions with the same denominator are ordered according to the size of the numerator. If, however, one wished to order the fractions according to the size of their numerical value, one would obtain no countable arrangement.

2. \gesperrt{The totality of all algebraic numbers is an enumerable set.}
